\section{Methodology}\label{methodology}

\subsection{\texorpdfstring{\textbf{3.1 Research Design}}{3.1 Research Design}}\label{research-design}

This study will use a quantitative, non-experimental design to model residential property values in Metro Cebu. The analysis will be predictive in that it will estimate prices from observed property characteristics, and prescriptive in that the results will later be organized for spatial decision support through GIS. The dependent variable will be property price, or where appropriate price per square meter, while the independent variables will include structural, geospatial, administrative, and macroeconomic features.

Three supervised learning approaches will be compared in order to balance interpretability and predictive performance, while also testing whether GIS-derived and administrative variables improve the model beyond structural features alone:

\begin{enumerate}
\item \textbf{Ordinary Least Squares (OLS) / Hedonic Regression}: An interpretable baseline grounded in Hedonic Pricing Theory (Rosen, 1974). Coefficients carry direct economic meaning (e.g., ``each additional bedroom adds PHP X to value'').
\item \textbf{Random Forest Regressor}: A tree-based ensemble that captures non-linear relationships and feature interactions without requiring explicit specification (Breiman, 2001).
\item \textbf{XGBoost Regressor}: A gradient boosting algorithm optimized for predictive accuracy on structured tabular data (Chen \& Guestrin, 2016). Empirically, XGBoost often performs well on datasets of similar scale (Nyanda et al., 2024).
\end{enumerate}

\begin{center}\rule{0.5\linewidth}{0.5pt}\end{center}

\subsection{\texorpdfstring{\textbf{3.2 Data Sources: The Hybrid Strategy}}{3.2 Data Sources: The Hybrid Strategy}}\label{data-sources-the-hybrid-strategy}

Because direct deed-of-sale transaction records are not publicly accessible, the dataset will be assembled from multiple sources that reflect different segments of the residential market. Foreclosure and acquired-asset listings will provide a more conservative segment of observed prices, while online listings will provide asking prices that reflect current market exposure more closely. The purpose of this hybrid strategy is not to claim direct observation of a single true market price, but to work with the most usable price signals available under existing data constraints.

\begin{center}
\begin{tabular}{p{0.26\textwidth} p{0.20\textwidth} p{0.22\textwidth} p{0.20\textwidth}}
\toprule
Source & Role & Volume & Nature \\
\midrule
\textbf{Bank Foreclosures} & Verified Floor Price & BDO, PNB, RCBC, Union Bank & Conservative / Distressed \\
\textbf{Pag-IBIG / Gov't Assets} & Verified Floor Price & HDMF acquired assets & Conservative / Distressed \\
\textbf{Online Listings} (Lamudi) & Market Ceiling Price & Target: 500+ & Asking / Speculative \\
\textbf{BIR Zonal Values} & Administrative Benchmark & Per barangay & Static / Regulatory \\
\textbf{BSP RPPI} & Time-Trend Control & Quarterly index & Macro / Cyclical \\
\textbf{Google Maps API} & Geocoding / Location Data & Per property & Geospatial / Dynamic \\
\textbf{OpenStreetMap (OSM)} & Amenity \& Land-Use Data & Per property radius & Geospatial / Open Data \\
\bottomrule
\end{tabular}
\end{center}

\subsubsection{\texorpdfstring{\textbf{3.2.1 Primary Data: Verified Floor Price (Foreclosed / Acquired Assets)}}{3.2.1 Primary Data: Verified Floor Price (Foreclosed / Acquired Assets)}}\label{primary-data-verified-floor-price-foreclosed-acquired-assets}

The floor price dataset will aggregate foreclosed and acquired property listings from multiple institutional sources to prevent single-source bias:

\begin{itemize}
\item \textbf{BDO Unibank}: 955 raw foreclosure entries nationwide (as of November 18, 2025).
\item \textbf{Pag-IBIG Fund (HDMF)}: Acquired assets representing the affordable housing segment.
\item \textbf{PNB, RCBC, and Union Bank}: Additional listings to broaden coverage.
\item \textbf{SSS / GSIS}: Government-acquired properties within Metro Cebu.
\end{itemize}

These distressed assets will be treated as a lower-bound segment of the observed market rather than as direct equivalents of open-market transaction prices. Using multiple institutional sources helps reduce the bias that might result if the dataset were drawn from only one lender or seller.

\textbf{Sample Data Structure: BDO Foreclosures (Raw vs.~Cleaned)}

\begin{center}
\begin{tabular}{p{0.28\textwidth} p{0.28\textwidth} p{0.30\textwidth}}
\toprule
Raw Column (BDO Excel) & Processed Feature & Description \\
\midrule
REGION & Region & Filtered to Region VII \\
CITY\_PROVINCE & City & Filtered to Metro Cebu LGUs \\
PROPERTY\_ADDRESS & Address & Geocoded string \\
LOT\_AREA & Lot Area & Numeric (sqm) \\
FLOOR\_AREA & Floor Area & Numeric (sqm) \\
MINIMUM\_BID\_PRICE & Actual Price & Numeric target variable (PHP) \\
PROPERTY\_DESCRIPTION & Bedrooms, Bathrooms & Parsed via regex \\
\bottomrule
\end{tabular}
\end{center}

\subsubsection{\texorpdfstring{\textbf{3.2.2 Secondary Data: Market Ceiling Price (Online Listings)}}{3.2.2 Secondary Data: Market Ceiling Price (Online Listings)}}\label{secondary-data-market-ceiling-price-online-listings}

To represent current asking prices in the residential market, the study also collects listings from public online platforms, primarily Lamudi. As Sousa et al. (2024) note, large listing datasets can capture pricing patterns that are often missing from sparse official records. For this reason, the study targets at least 500 Metro Cebu residential listings.

\textbf{Sample Data Structure: Lamudi Listings (Raw vs.~Cleaned)}

\begin{center}
\begin{tabular}{p{0.26\textwidth} p{0.26\textwidth} p{0.34\textwidth}}
\toprule
Raw Field (Web Scrape) & Processed Feature & Description \\
\midrule
Location & Barangay, City & Parsed location string \\
Price & Actual Price & Target variable (PHP), cleaned of currency symbols \\
Bedrooms & Bedrooms & Numeric extract \\
Bathrooms & Bathrooms & Numeric extract \\
Floor area & Floor Area & Numeric extract (sqm) \\
Land Size & Lot Area & Numeric extract (sqm) \\
Latitude / Longitude & Latitude, Longitude & Direct coordinate mapping \\
\bottomrule
\end{tabular}
\end{center}

\subsubsection{\texorpdfstring{\textbf{3.2.3 Administrative and Macroeconomic Data}}{3.2.3 Administrative and Macroeconomic Data}}\label{administrative-and-macroeconomic-data}

\begin{itemize}
\item \textbf{BIR Zonal Values}: Official zonal values per barangay, to be used both as a model feature and to calculate the ``Valuation Gap'' (Market Price - Zonal Value).
\item \textbf{BSP RPPI}: Quarterly Residential Real Estate Price Index for Areas Outside NCR (AONCR), controlling for time-trend effects (inflation/market cycle).
\end{itemize}

\subsubsection{\texorpdfstring{\textbf{3.2.4 Geospatial Data Sources}}{3.2.4 Geospatial Data Sources}}\label{geospatial-data-sources}

\begin{itemize}
\item \textbf{Google Maps Geocoding API}: Will convert property addresses into precise latitude/longitude coordinates, enabling all subsequent spatial analyses. Chosen for its superior address disambiguation in Philippine contexts.
\item \textbf{OpenStreetMap (OSM)}: Will provide open geospatial data for amenity density analysis. Will be queried via the osmnx Python library (Boeing, 2017) to retrieve counts of schools, hospitals, commercial establishments, restaurants, and public transport stops within defined radii of each property.
\end{itemize}

\subsubsection{\texorpdfstring{\textbf{3.2.5 Target Variable}}{3.2.5 Target Variable}}\label{target-variable}

Within this hybrid dataset, foreclosure prices and online listing prices will not be treated as identical market signals. A source indicator will therefore be included so the model can account for systematic level differences between distressed and market-facing listings. This will allow the analysis to estimate price relationships across both segments without assuming that they represent the same selling condition.

\textbf{Final Feature Matrix Schema (Pre-Modeling)}

\begin{center}
\begin{tabular}{p{0.17\textwidth} p{0.55\textwidth} p{0.18\textwidth}}
\toprule
Feature Group & Variables & Type \\
\midrule
\textbf{Identifiers} & ID, Source (BDO/Lamudi) & Categorical \\
\textbf{Structural} & Lot Area, Floor Area, Bedrooms, Bathrooms, Property Type & Numeric / Categorical \\
\textbf{Locational} & Latitude, Longitude, Barangay & Numeric / Categorical \\
\textbf{Geospatial} & Dist\_CBD, Dist\_Airport, Dist\_CBRT & Numeric (Meters) \\
\textbf{Amenity} & OSM\_Amenity\_Score & Numeric (Count/Index) \\
\textbf{Economic} & BIR\_Zonal\_Value, Spatial\_Lag\_Mean & Numeric (PHP) \\
\textbf{Target} & Actual Price, Log\_Price & Numeric (PHP) \\
\bottomrule
\end{tabular}
\end{center}

\subsubsection{\texorpdfstring{\textbf{3.2.6 Validation Layer: Human-in-the-Loop}}{3.2.6 Validation Layer: Human-in-the-Loop}}\label{validation-layer-human-in-the-loop}

Licensed real estate brokers from the CPRE network will be included as a validation layer, not to replace statistical evaluation, but to check whether the model outputs remain plausible within local market practice. Their role will be limited to two tasks:

\begin{enumerate}
\item \textbf{Sanity Check}: Reviewing SHAP-derived value driver rankings against practitioner knowledge.
\item \textbf{Outlier Review}: Examining cases with high prediction error to distinguish likely data issues from genuine market anomalies.
\end{enumerate}

\begin{center}\rule{0.5\linewidth}{0.5pt}\end{center}

\subsection{\texorpdfstring{\textbf{3.3 Data Pipeline}}{3.3 Data Pipeline}}\label{data-pipeline}

The data preparation process proceeds in five stages.

\begin{enumerate}
\item \textbf{Ingestion}: We will ingest BDO Excel files via Pandas and will collect Lamudi listings using web scraping logic.
\item \textbf{Filtering}: The dataset will be restricted to residential properties within Metro Cebu (Cebu City, Mandaue, Lapu-Lapu, Talisay, Minglanilla, and Consolacion).
\item \textbf{Regex Parsing and Cleaning}:
\begin{itemize}
\item \emph{BDO Data}: The Property Description string often bundles features (e.g., ``3BR 2TB''). We will apply regex patterns to extract Bedrooms and Bathrooms. Missing values in lot/floor area will be imputed based on the median of similar property types in the same barangay.
\item \emph{Lamudi Data}: Scraped fields often contain varying text formats (e.g., ``PHP 5,000,000'' or ``Contact agent for price''). We will clean currency symbols and commas, dropping rows without explicit numerical prices.
\end{itemize}
\item \textbf{Geocoding}: We will batch-process addresses through the Google Maps Geocoding API, standardizing barangay names and securing precise coordinates.
\item \textbf{GIS Augmentation}: From these geocoded coordinates, we will engineer geospatial features:
\begin{itemize}
\item Proximity metrics (Haversine distances to key economic nodes).
\item Amenity scores (OSM-derived counts within a 1 km radius).
\item Spatial lag (mean price of neighboring properties within a 1 km radius).
\end{itemize}
\end{enumerate}

\textbf{Tools}: Python (Pandas, Scikit-learn, XGBoost, osmnx), Google Maps API for geocoding, QGIS for spatial visualization.

\begin{center}\rule{0.5\linewidth}{0.5pt}\end{center}

\subsection{\texorpdfstring{\textbf{3.4 Feature Engineering}}{3.4 Feature Engineering}}\label{feature-engineering}

\begin{center}
\begin{tabular}{p{0.18\textwidth} p{0.58\textwidth} p{0.16\textwidth}}
\toprule
Category & Features & Source \\
\midrule
\textbf{Structural} & Lot Area, Floor Area, Bedrooms, Bathrooms, Parking, Property Type & BDO / Lamudi \\
\textbf{Locational} & Barangay, Latitude/Longitude & Google Maps API \\
\textbf{Geospatial} & Proximity to Ayala, IT Park, SM Seaside, Airport, \textbf{CBRT stations} (Haversine) & Geocoding + GIS \\
\textbf{Amenity Score} & Count of schools, hospitals, commercial centers, transit stops within 1 km radius & OSM / osmnx \\
\textbf{Spatial Lag} & Mean price of neighboring properties within defined radius & Computed from data \\
\textbf{Administrative} & BIR Zonal Value (per barangay) & BIR schedules \\
\textbf{Macro} & BSP RPPI quarterly index & BSP data \\
\textbf{Data Source} & Source indicator (BDO vs.~Lamudi) & Engineered \\
\bottomrule
\end{tabular}
\end{center}

\textbf{Engineered variables}: Price per sqm, Valuation Gap (Price - Zonal Value), Log(Price).

\subsubsection{\texorpdfstring{\textbf{3.4.1 Geospatial Feature Engineering}}{3.4.1 Geospatial Feature Engineering}}\label{geospatial-feature-engineering}

Geospatial feature construction is the part of the method that most directly connects the valuation model to the urban structure of Metro Cebu. Rather than treating location as a background descriptor, this stage translates accessibility, amenity access, and neighborhood context into variables that can enter the model explicitly.

\begin{enumerate}
\item \textbf{Geocoding (Google Maps API)}: Each property address will be geocoded to obtain precise latitude/longitude coordinates. The Google Maps Geocoding API has been chosen for its superior handling of Philippine address formats, which often include barangay names, landmarks, or informal location descriptors.
\item \textbf{Proximity Features (Haversine Formula)}: For each geocoded property, the Haversine formula will compute great-circle distances to key economic and infrastructure nodes:
\end{enumerate}

\begin{itemize}
\item Ayala Center Cebu (primary CBD)
\begin{itemize}
\item Cebu IT Park (employment hub)
\item SM Seaside City (commercial center)
\item Mactan-Cebu International Airport
\item Planned Cebu Bus Rapid Transit (CBRT) station locations
\end{itemize}
\end{itemize}

\begin{enumerate}
\setcounter{enumi}{2}
\item \textbf{Custom Value Driver Scoring Model (OSM via osmnx)}: Rather than relying only on separate distance measures, the study will also construct an amenity score from OpenStreetMap data. Using the osmnx library, points of interest within a 1 kilometer network radius of each property will be queried and grouped into categories relevant to residential valuation. A 1 km radius is used because it roughly corresponds to a 10 to 15 minute walkable catchment in urban settings. The resulting score is designed to reflect neighborhood service access rather than simple amenity counts alone.
\begin{itemize}
\item Educational institutions (schools, universities): Standard weight
\item Healthcare facilities (hospitals, clinics): High weight
\item Commercial establishments (malls, markets): Medium weight
\item Public transport stops (jeepney routes, bus stops): High weight
\end{itemize}
\emph{Note: Initial weights will be drawn from the literature and then checked during exploratory analysis so that no single amenity category dominates the index without empirical basis.}
\item \textbf{Spatial Lag Variable}: To capture neighborhood price effects (spatial autocorrelation), we will compute the mean actual price of all other properties within a 1 kilometer radius of the target property. This will operationalize Tobler's First Law---that near things are more related than distant things---directly into our non-spatial ML algorithms.
\end{enumerate}

\begin{center}\rule{0.5\linewidth}{0.5pt}\end{center}

\subsection{\texorpdfstring{\textbf{3.5 Pre-processing}}{3.5 Pre-processing}}\label{pre-processing}

\begin{center}
\begin{tabular}{p{0.24\textwidth} p{0.28\textwidth} p{0.34\textwidth}}
\toprule
Step & Method & Rationale \\
\midrule
\textbf{Outlier Detection} & IQR (Interquartile Range) method & Statistically principled; no arbitrary cutoffs \\
\textbf{Log Transformation} & ln(Price) as target variable & Normalizes the right-skewed price distribution \\
\textbf{Missing Values} & Barangay-level median imputation & Preserves local spatial context \\
\textbf{Encoding} & One-Hot Encoding (Property Type, Barangay) & Required for regression and tree-based models \\
\bottomrule
\end{tabular}
\end{center}

\begin{center}\rule{0.5\linewidth}{0.5pt}\end{center}

\subsection{\texorpdfstring{\textbf{3.6 Modeling Strategy}}{3.6 Modeling Strategy}}\label{modeling-strategy}

\subsubsection{\texorpdfstring{\textbf{3.6.1 The Three Models}}{3.6.1 The Three Models}}\label{the-three-models}

\begin{center}
\begin{tabular}{p{0.05\textwidth} p{0.28\textwidth} p{0.38\textwidth} p{0.18\textwidth}}
\toprule
\# & Model & Strength & Weakness \\
\midrule
1 & \textbf{Hedonic Regression (OLS)} & Interpretable; coefficients have economic meaning & Assumes linearity \\
2 & \textbf{Random Forest} & Handles non-linearities; robust to overfitting & Less interpretable \\
3 & \textbf{XGBoost} & Best predictive performance on tabular data & Hyperparameter-sensitive \\
\bottomrule
\end{tabular}
\end{center}

\subsubsection{\texorpdfstring{\textbf{3.6.2 Hedonic Equation}}{3.6.2 Hedonic Equation}}\label{hedonic-equation}

The hedonic regression will take the following log-linear form:

\[
\begin{aligned}
\ln(Price) ={}& \alpha + \beta_1 \ln(Area) + \beta_2(BR) + \beta_3(Dist_{CBD}) \\
& + \beta_4(ZonalValue) + \beta_5(AmenityScore) + \beta_6(SpatialLag) + \epsilon
\end{aligned}
\]

In this form, the hedonic model retains its interpretive structure while allowing administrative and spatial variables to enter directly into the price equation. This makes it possible to compare a more conventional valuation framework with the tree-based models without stripping out the locational features that are central to the study.

\subsubsection{\texorpdfstring{\textbf{3.6.3 Hyperparameter Tuning}}{3.6.3 Hyperparameter Tuning}}\label{hyperparameter-tuning}

For the ML models (Random Forest and XGBoost), hyperparameters will be optimized via \textbf{GridSearchCV} with \textbf{K-Fold Cross Validation} (K = 5 or 10, depending on final sample size). This will avoid overfitting while maximizing generalization performance.

\begin{center}\rule{0.5\linewidth}{0.5pt}\end{center}

\subsection{\texorpdfstring{\textbf{3.7 Evaluation and Explainability}}{3.7 Evaluation and Explainability}}\label{evaluation-and-explainability}

\subsubsection{\texorpdfstring{\textbf{3.7.1 Performance Metrics}}{3.7.1 Performance Metrics}}\label{performance-metrics}

\begin{center}
\begin{tabular}{p{0.14\textwidth} p{0.28\textwidth} p{0.44\textwidth}}
\toprule
Metric & Description & Purpose \\
\midrule
\textbf{MAPE} & Mean Absolute Percentage Error & Primary metric; enables cross-study comparison \\
\textbf{R\textsuperscript{2}} & Coefficient of Determination & Proportion of variance explained \\
\textbf{MAE} & Mean Absolute Error (in PHP) & Business-interpretable error \\
\textbf{RMSE} & Root Mean Square Error & Penalizes large errors \\
\bottomrule
\end{tabular}
\end{center}

\subsubsection{\texorpdfstring{\textbf{3.7.2 Benchmark Targets}}{3.7.2 Benchmark Targets}}\label{benchmark-targets}

\begin{center}
\begin{tabular}{p{0.22\textwidth} p{0.20\textwidth} p{0.08\textwidth} p{0.36\textwidth}}
\toprule
Study & Context & MAPE & Our Target \\
\midrule
Ramolete et al. (2023) & Philippines, larger dataset & 10--21\% & \textless{} 25\% (they had larger, cleaner data) \\
Nyanda et al. (2024) & Tanzania, $n \approx 954$ & 48\% & Beat 48\% (same sample size, but we add GIS features) \\
\bottomrule
\end{tabular}
\end{center}

\subsubsection{\texorpdfstring{\textbf{3.7.3 SHAP Explainability}}{3.7.3 SHAP Explainability}}\label{shap-explainability}

SHAP will be used here to interpret model outputs at both the dataset and property level, which is important if the results are to remain reviewable in valuation practice and consistent with IVS 2025 transparency requirements.

\begin{itemize}
\item \textbf{Global SHAP (Summary Plots)}: Will identify which value drivers affect Metro Cebu property prices most across the entire dataset. This will answer RQ1 (``What value drivers significantly influence property prices in Metro Cebu?'').
\item \textbf{Local SHAP (Force Plots)}: Will explain individual predictions. For example: \emph{``This Lahug condo is valued at PHP X: +PHP 1.2M due to IT Park proximity, -PHP 300K due to small floor area, +PHP 200K due to high amenity score.''}
\end{itemize}

This helps keep the model interpretable enough to function as a support tool rather than a purely opaque prediction system. In practical terms, the model is evaluated not only by how accurately it predicts price, but also by whether its outputs can still be examined and discussed in terms that practitioners can understand.

\begin{center}\rule{0.5\linewidth}{0.5pt}\end{center}

\subsection{\texorpdfstring{\textbf{3.8 Deliverables}}{3.8 Deliverables}}\label{deliverables}

\subsubsection{\texorpdfstring{\textbf{3.8.1 QGIS Interactive Map}}{3.8.1 QGIS Interactive Map}}\label{qgis-interactive-map}

The main applied output of the study is an interactive QGIS project that organizes model results spatially. Rather than presenting the results only as tables or summary statistics, the map is intended to show how predicted values, valuation gaps, and locational patterns are distributed across Metro Cebu.

\begin{enumerate}
\item \textbf{Property Valuations}: Point vectors representing individual geocoded properties. They will be color-coded based on the model's prediction error (actual vs.~predicted), allowing users to visually identify undervalued anomalies or overvalued clusters.
\item \textbf{Valuation Gap Heatmap}: A raster heatmap visualizing the divergence between the ML model predictions and the official BIR Zonal Values. Hotspots will indicate areas where official valuations significantly lag market realities.
\item \textbf{CBRT \& Infrastructure Overlays}: Line segments denoting the planned CBRT route with 500m and 1km buffer zones. This will allow users to visualize how upcoming public transit infrastructure intersects with current market valuations.
\item \textbf{Value Driver Contours}: ISO-chrones or distance contours measuring proximity to the CBD (Ayala Center) or IT Park.
\end{enumerate}

The purpose of this output is to make the model easier to inspect geographically, especially for brokers, investors, and local government users who need to compare values across locations rather than only at the level of individual records.

\subsubsection{\texorpdfstring{\textbf{3.8.2 Streamlit Web Application}}{3.8.2 Streamlit Web Application}}\label{streamlit-web-application}

A Streamlit web application will complement the QGIS map by providing an interactive interface for testing individual property inputs against the trained model. Users will be able to enter structural property features, select a location, and receive a predicted price together with a SHAP-based explanation of the factors contributing to that estimate. While the QGIS output emphasizes spatial comparison across Metro Cebu, the Streamlit application provides a more direct way to inspect individual predictions and their underlying value drivers.

\begin{center}\rule{0.5\linewidth}{0.5pt}\end{center}

\subsection{\texorpdfstring{\textbf{3.9 Timeline and Milestones}}{3.9 Timeline and Milestones}}\label{timeline-and-milestones}

\begin{center}
\begin{tabular}{p{0.22\textwidth} p{0.48\textwidth} p{0.18\textwidth}}
\toprule
Phase & Activity & Timeline \\
\midrule
\textbf{1. Data} & Lamudi scraping + BDO cleaning & Feb 18 -- Feb 28 \\
\textbf{2. Proposal} & Panel Presentation & Feb 21 \\
\textbf{3. Build} & Geocoding + GIS feature engineering (Google Maps, OSM) & Mar 1 -- Mar 14 \\
 & Model training + Hyperparameter tuning & Mar 15 -- Mar 28 \\
\textbf{4. Colloquium} & Research updates presentation & Mar 28 \\
\textbf{5. Evaluate} & SHAP analysis + QGIS map + Broker validation & Apr 1 -- Apr 18 \\
\textbf{6. Write} & Draft final paper (Chapters 4--10) & Apr 18 -- May 2 \\
\textbf{7. Defend} & Final Research Paper presentation & May 9 \\
 & Final Submission & May 23 \\
\bottomrule
\end{tabular}
\end{center}
